\chapter{结论与展望}

本文面向严格零样本跨被试 fMRI 到图像重建任务，围绕“共享视觉语义”与“个体神经差异”在有限观测数据下高度耦合这一核心瓶颈，提出了统一框架 MINDMELD。与传统方法将跨被试差异视为噪声或仅通过后验适配进行补偿不同，本文从建模机制上同时刻画被试相关变化与被试共享语义，通过“被试感知神经响应预测模块”与“被试不变语义解码模块”的协同设计，将原本阻碍跨被试泛化的个体差异，转化为可被利用的结构化监督信息，从而为无需目标被试校准的视觉重建提供了新的解决思路。

具体而言，本文首先构建了被试感知神经响应预测模块，在统一语义输入下显式建模不同被试之间的响应变化规律。该模块不仅增强了模型对神经信号分布的刻画能力，也在训练过程中有效扩展了原始神经观测流形，使模型能够在更丰富的跨被试变化模式上学习稳定的映射关系。进一步地，本文设计了被试不变语义解码模块，通过语义 token 学习、对抗解耦与跨重建约束，逐步剥离神经响应中的个体特异性成分，提取能够跨被试共享的语义核心表示。两部分相互配合，形成“生成监督增强表征学习、语义约束反哺跨被试泛化”的闭环优化过程，使模型既具备处理神经多样性的能力，又能够保持语义表示的一致性与可迁移性。

在实验验证方面，本文基于 Natural Scenes Dataset 构建了严格的零样本跨被试评估协议，并在多个目标被试上系统比较了所提方法与当前代表性方法的表现。实验结果表明，MINDMELD 在零样本设置下取得了优异的综合性能：一方面，模型在 SSIM、Alex(2)、Inception、CLIP 以及感知距离等多个关键指标上表现突出，说明其能够同时提升重建图像的低层结构保真度与高层语义一致性；另一方面，定性结果也进一步表明，MINDMELD 在主体语义恢复、场景整体布局保持以及多物体复杂场景的稳定表达方面均具有明显优势。这说明本文方法不仅在数值评价上取得提升，更在实际视觉可感知质量上展现出较强的泛化能力。

除整体性能外，本文还通过消融实验与表示分析，对模型有效性进行了更细致的机制验证。消融研究表明，被试感知神经响应预测模块能够为跨被试学习提供更充分的结构先验，语义解耦结构有助于稳定提取共享语义表示，而对抗约束与一致性约束则共同保证了解耦过程的充分性与语义空间的稳定性。进一步的表示分析显示，模型在没有引入显式解剖先验或层级监督的情况下，自发学习到了与视觉皮层层级组织相一致的特征偏好模式；同时，语义表示在潜在空间中呈现出“被试混合、语义分离”的良好结构，表明模型成功构建了兼具被试不变性与任务判别性的语义嵌入。这些结果从机理层面解释了 MINDMELD 能够实现零样本跨被试泛化的原因，也体现了本文方法在可解释性方面的优势。

总体而言，本文工作的主要结论可以概括为以下三点。第一，在零样本跨被试视觉重建任务中，单纯依赖被试对齐或目标被试微调并不足以从根本上解决泛化问题，关键在于显式区分神经响应中的共享语义成分与个体差异成分。第二，通过将被试感知的响应生成与被试不变的语义建模进行联合优化，可以将神经多样性由干扰因素转变为有益监督，从而显著提升模型对未见被试的适应能力。第三，在语义层施加跨被试一致性约束，是构建可迁移、可解释且适用于生成模型驱动重建任务的潜在表示空间的一条有效路径。上述结论表明，围绕“合成神经监督”与“跨被试语义约束”构建闭环学习框架，为实现高鲁棒性、高可扩展性的免校准视觉脑机接口提供了切实可行的方法基础。

当然，本文工作仍存在一定局限性。首先，当前实验主要基于 NSD 数据集开展，尽管该数据集规模较大、标注充分，但其任务范式、采样条件和被试数量仍然有限，模型在更广泛数据采集条件下的适用性还有待进一步验证。其次，本文重点关注的是语义层面对齐与图像重建效果，对于时间动态信息、认知状态变化以及更细粒度脑区功能差异的建模仍不充分。再次，当前框架主要面向静态自然图像重建，尚未扩展到视频、想象、记忆回忆等更复杂的神经解码场景。此外，尽管模型在解释性方面已表现出较好潜力，但潜在表示与具体神经机制之间的对应关系仍需结合更多神经科学分析工具做进一步确认。

未来工作可从以下几个方向展开。其一，可将所提出的跨被试语义建模思想推广到更大规模、多中心、跨采集设备的数据上，进一步提升模型对真实场景中数据偏移的适应能力。其二，可结合时序建模方法，将当前框架扩展到动态视觉刺激或连续认知过程的解码任务中，从而提升模型对复杂脑活动过程的刻画能力。其三，可进一步探索与更强生成模型的结合方式，在保持语义稳定性的同时提升细节还原能力和图像自然度。其四，可引入更细粒度的神经科学先验，对潜在表示与皮层功能组织之间的关系进行联合分析，以增强模型的生物可解释性与科学发现价值。其五，从应用角度看，若能进一步降低计算成本并提高模型推理稳定性，本文方法有望为面向真实用户的免校准视觉脑机接口系统提供支撑，并推动跨被试神经解码技术向更高实用性和可部署性发展。

综上所述，本文提出的 MINDMELD 为零样本跨被试 fMRI 到图像重建问题提供了一种统一、有效且具有可解释性的解决方案。通过打通“神经响应扩展—语义解耦—跨被试重建”之间的优化链路，本文不仅提升了未见被试上的重建质量，也从表示学习与神经机制两个层面加深了对跨被试语义建模问题的理解。我们相信，随着更丰富数据、更强生成模型与更深入神经科学分析方法的不断结合，基于共享语义空间的跨被试脑解码框架将在未来展现出更大的理论价值与应用潜力。

\chapterbib
